%% ============================================================
%% Formatting adapted to match the Springer Nature single-column
%% structure as seen in akshat_paper.pdf
%% ============================================================
\documentclass[12pt]{article}

%% --- Packages ---
\usepackage[T1]{fontenc}
\usepackage[utf8]{inputenc}
\usepackage{geometry}
\geometry{a4paper, margin=1in}
\usepackage{authblk}
\usepackage{graphicx}
\usepackage{amsmath}
\usepackage{amssymb}
\usepackage{booktabs}
\usepackage{array}
\usepackage{multirow}
\usepackage{enumitem}
\usepackage{float}
\usepackage{url}
\usepackage{xcolor}
\usepackage{cite}
\usepackage{textcomp}
\usepackage{subcaption}
\usepackage{wasysym}
\newcommand{\partmark}{\LEFTcircle}
\usepackage{algorithm}
\usepackage{algpseudocode}
\usepackage{hyperref}

%% Graphics path
\graphicspath{{./}{../report/}{../}}

%% hyperref setup
\hypersetup{
  colorlinks=true,
  linkcolor=blue!60!black,
  citecolor=green!50!black,
  urlcolor=blue!70!black
}

\renewcommand\Authfont{\normalsize}
\renewcommand\Affilfont{\itshape\small}

%% ============================================================
\begin{document}

\title{\Large\bfseries Agentic AI-Assisted Autonomous Orchestration for 5G Network
Slicing: A LangGraph Multi-Agent Framework with Chain-of-Thought
Reasoning and Wrong-Lever Avoidance}

\author[1]{Simran R. Kalkeri}
\author[1]{Amaanali D. Doddamani}
\author[1]{Sparsh S. Naik}
\author[1]{Rishi P. Kulkarni}
\author[1]{Narayan D. G.*}

\affil[1]{School of Computer Science and Engineering, KLE Technological University,
Hubballi -- 580031, Karnataka, India.}
\affil[ ]{\textit{*Corresponding author. E-mail: narayan.dg@kletech.ac.in;}}
\affil[ ]{\textit{Contributing authors: simran.kalkeri@kletech.ac.in; amaanali.doddamani@kletech.ac.in; sparsh.naik@kletech.ac.in; rishi.kulkarni@kletech.ac.in;}}

\date{}

\maketitle

%% ============================================================
%% ABSTRACT
%% ============================================================
\begin{abstract}
\noindent Fifth-generation (5G) network slicing enables the concurrent operation
of enhanced Mobile Broadband (eMBB), Ultra-Reliable Low-Latency
Communication (URLLC), and massive Machine-Type Communication (mMTC)
services on shared physical infrastructure. Static threshold-based
controllers, while operationally simple, fail under dynamic traffic
conditions: they apply control actions without causal reasoning, waste
bandwidth by throttling idle slices, and retain no operational memory
across control cycles. This paper presents an autonomous, cloud-native
agentic orchestration framework for 5G network slice management. The
system integrates a LangGraph-based multi-agent pipeline with a locally
hosted Large Language Model (LLM)---specifically \texttt{llama3.2:3b}
via Ollama---to perform structured Chain-of-Thought (CoT) reasoning for
slice QoS enforcement. The framework contributes a five-VM cloud-native
5G testbed combining Open5GS, UERANSIM, and Kubernetes-hosted per-slice
User Plane Functions (UPFs); a decoupled multi-agent architecture spanning
Monitoring, State, Planning, Validation, and Execution agents; a mandatory
JSON CoT schema requiring root-cause assessment and lever-validity scoring
before any control action; a Wrong-Lever Avoidance (WLA) mechanism using
the novel $\rho_\text{eMBB}$ load-fraction metric to suppress spurious
throttling; and an episodic memory subsystem enabling memory-assisted
adaptive decision-making. Experimental evaluation against a rule-based
baseline demonstrates that the agentic orchestrator improves URLLC SLA
compliance from 87.1\% to 96.3\%, halves mean recovery time from 8.4\,s
to 4.2\,s, reduces unnecessary eMBB throughput loss by 72\%, and maintains
decision overhead below 380\,ms---well within the 2-second macroscopic
control window.
\end{abstract}

\medskip
\noindent\textbf{Keywords:} 5G network slicing, agentic AI, LLM orchestration, LangGraph, QoS enforcement, URLLC, eMBB, chain-of-thought reasoning, cloud-native, Kubernetes, Open5GS, multi-agent systems
\vspace{1em}

%% ============================================================
%% 1. INTRODUCTION
%% ============================================================
\section{Introduction}
\label{sec:introduction}

Fifth Generation (5G) mobile networks introduce network slicing as a
defining architectural capability, enabling multiple logical network
instances---each with distinct performance characteristics---to coexist
on shared physical infrastructure~\cite{tsai2021upf}. Network slicing
permits concurrent support of enhanced Mobile Broadband (eMBB),
Ultra-Reliable Low-Latency Communication (URLLC), and massive
Machine-Type Communication (mMTC) services, each with fundamentally
different Quality-of-Service (QoS) requirements. eMBB demands high
throughput and tolerates moderate latency; URLLC requires end-to-end
latencies below 20\,ms with very high reliability; and mMTC targets
energy-efficient transmission for large-scale IoT device populations.

Realizing these guarantees in practice is challenging. At the user plane,
traffic from different slices may contend for shared buffering and
forwarding resources during congestion. Static resource allocations fail
to track dynamic traffic demands, while reactive threshold-based
controllers lack the ability to reason about the root cause of degradation
or to anticipate violations before they occur~\cite{zhou2021predict}.
Autonomous orchestration, enabled by Large Language Models (LLMs) and
multi-agent frameworks, offers a qualitatively different control paradigm.
Rather than pattern-matching against hardcoded thresholds, an LLM-based
planning agent performs causal reasoning: identifying which network entity
is responsible for congestion, validating whether the proposed control
action addresses that cause, and consulting prior operational history
stored in an agent memory to refine its decisions~\cite{park2024llm,openai2023cot}.
This paper implements, deploys, and evaluates such a system in a
cloud-native 5G testbed.

\medskip\noindent\textbf{Motivation.}\quad
The evolution from rule-based to autonomous orchestration is driven by
three fundamental limitations observed in practical 5G slice management
systems. \textbf{First}, rule-based controllers are structurally blind to
causation: a threshold-based system that detects elevated URLLC
Round-Trip Time (RTT) will throttle eMBB bandwidth regardless of whether
eMBB traffic is actually active. When eMBB load is low, this action is
unnecessary and harmful. An agent capable of reasoning about relative
load---quantified by the $\rho_\text{eMBB}$ metric introduced in this
work---can correctly distinguish idle from bursting eMBB and decline to
act when throttling is ineffective. \textbf{Second}, reactive control
cannot prevent violations; it can only respond to them. A URLLC SLA
breach lasting even 200\,ms may disrupt a latency-critical industrial
automation session, making pre-emptive, trend-aware intervention
essential~\cite{tran2025closed}. \textbf{Third}, accumulated operational
experience is wasted in memoryless systems: an agent that stores the
outcome of prior actions can converge to more accurate control policies
over time without requiring offline training. These three limitations
collectively motivate the agentic, causally-grounded, memory-augmented
design proposed in this work.

\medskip\noindent\textbf{Contributions.}\quad
This paper makes five interconnected contributions to autonomous 5G slice
orchestration. First (C1), we build a five-VM cloud-native 5G research
testbed integrating Open5GS, UERANSIM, Kubernetes-hosted per-slice UPFs,
and Prometheus\slash Grafana monitoring, providing the live experimental
substrate on which all claims are validated. Second (C2), we design a
LangGraph-based multi-agent orchestration pipeline with a monitoring
thread that is fully decoupled from LLM inference, ensuring continuous
3-second metric freshness regardless of planning latency. Third (C3), we
introduce a mandatory Chain-of-Thought JSON reasoning schema that forces
the LLM to produce a root-cause attribution and lever-validity score before
selecting any control action, making every decision fully auditable.
Fourth (C4), we propose a Wrong-Lever Avoidance (WLA) mechanism anchored
on the dimensionless $\rho_\text{eMBB}$ load-fraction metric, which
deterministically suppresses spurious throttling of idle slices and is the
primary driver of throughput preservation gains. Fifth (C5), we contribute
an episodic memory subsystem that encodes recent action outcomes into the
LLM planning context, enabling the agent to adapt its strategy based on
live operational history without any offline training.

The remainder of this paper is organized as follows.
Section~\ref{sec:background} reviews related work and identifies the
research gap. Section~\ref{sec:architecture} describes the testbed, the
seven-layer architecture, and the full agentic framework.
Section~\ref{sec:results} presents and analyses experimental results
across all five contributions. Section~\ref{sec:conclusion} concludes
and outlines future directions.

%% ============================================================
%% 2. BACKGROUND WORK
%% ============================================================
\section{Background Work}
\label{sec:background}

Research in autonomous 5G slice management spans three broadly overlapping
themes: dynamic orchestration of core-network functions, AI and
reinforcement-learning driven resource allocation, and LLM-assisted
intent-based management. We review representative work in each category
before identifying the gap that motivates our design.

\subsection{Dynamic Network Slice Orchestration}

Grings et al.~\cite{grings2022dynamic} propose a full dynamic
orchestration framework for 5G core network slicing over a cloud-native
Kubernetes platform. Unlike prior approaches that support only partial
orchestration through resource scaling, their solution introduces a
Kubernetes-integrated controller capable of dynamically reconfiguring 5G
core functions---including NSSF, AMF, SMF, and UPF---without terminating
active slices. Evaluation on a free5GC-based testbed showed a 47.5\%
reduction in slice reconfiguration requests compared to partial
orchestration. Novanana et al.~\cite{novanana2024mano} investigate
Kubernetes-based MANO for provisioning coexisting eMBB and URLLC services
using free5GC, UERANSIM, and OpenStack. Comparing single-UPF and dedicated
SMF+UPF scenarios, they demonstrate that dedicated isolation delivers more
stable throughput (20.2--22.1\,Gbps) and fewer retransmissions---a
finding that directly motivates our per-slice UPF deployment strategy.

\subsection{AI and Reinforcement Learning for Slice Management}

Reddy~\cite{reddy2025drl} proposes a Deep Reinforcement Learning
(DRL)-based Kubernetes autoscaling framework achieving 71.7\% faster
slice deployment, 35\% higher resource utilization, and 22\% lower
latency versus threshold-based autoscaling. Sulaiman
et al.~\cite{sulaiman2024microopt} present MicroOpt, combining DNN-based
traffic prediction with gradient optimization to achieve up to 21.9\%
lower resource consumption than state-of-the-art methods. Saha
et al.~\cite{saha2022monarch} design MonArch, a scalable monitoring
architecture for cloud-native 5G deployments, demonstrating that a
5-second Prometheus monitoring interval provides an effective
accuracy-overhead balance. While these approaches yield strong results
in simulation or partial deployments, they require offline training data
and cannot express the causal reasoning necessary to avoid wrong-lever
interventions.

\subsection{LLM-Assisted Network Management}

Park et al.~\cite{park2024llm} demonstrate LLM-based network management
spanning intent specification through to configuration generation.
Dandoush et al.~\cite{dandoush2024llm} propose a conceptual LLM-powered
framework integrating multi-agent systems with ETSI MANO and 3GPP slicing
architectures for natural-language intent translation, though without
experimental validation. Bandara et al.~\cite{bandara2026native} present
an agentic AI control plane for 6G slice orchestration using the Model
Context Protocol (MCP), demonstrating successful LLM fine-tuning and
reliable tool invocation on an Open5GS, Ericsson RAN, and Kubernetes
testbed. Tran et al.~\cite{tran2025closed} propose PCLANSA, an
LSTM-based closed-loop service assurance framework for 5G/B5G slicing
achieving resource savings of 54.85\% for eMBB and 57.1\% for URLLC
versus static provisioning.

\subsection{Research Gap}

Existing approaches either rely on deterministic rule engines that lack
causal reasoning, offline-trained DRL agents that require extensive
simulation data, or conceptual LLM frameworks that have not been validated
on live networks. None provides an LLM-based, causally-grounded,
memory-augmented orchestration system evaluated on a live cloud-native 5G
testbed with real user-plane isolation simultaneously across three slice
types. Our work fills this gap by combining live testbed operation,
structured causal reasoning, deterministic safety gating, and episodic
memory into a single coherent framework.

%% ============================================================
%% 3. SYSTEM ARCHITECTURE
%% ============================================================
\section{System Architecture}
\label{sec:architecture}

This section describes the complete system from physical infrastructure
through to the agentic decision loop. The design philosophy is layered:
the lower layers provide the 5G substrate and traffic isolation machinery,
while the upper layers implement the autonomous control intelligence.
Figure~\ref{fig:sysarch} shows the overall architecture of the proposed
system.

\begin{figure}[H]
\centering
\rule{\textwidth}{6cm} \\ \vspace{0.5em}\centering\textbf{[Image Placeholder: system\_archie\_3.png]}
\caption{End-to-end architecture of the proposed Agentic 5G Slice
Orchestrator. User traffic generated by UERANSIM traverses the Open5GS
5G core, per-slice UPFs, and Linux \texttt{tc} shaping before reaching
MEC applications. The agentic orchestration layer on kubemaster closes
the control loop via Prometheus, SSH, and the Kubernetes API.}
\label{fig:sysarch}
\end{figure}

\subsection{Testbed Infrastructure}

The testbed is deployed on a single physical host running VMware, hosting
five virtual machines sharing a flat Layer-2 network
($192.168.49.0/24$). Table~\ref{tab:vms} summarizes the configuration
of each VM. The physical host provides Intel Xeon Gold 5418Y processors,
and the worker node (kube) is deliberately over-provisioned with
16\,vCPUs and 62\,GiB RAM to accommodate simultaneous UPF pods, MEC
application pods, and heavy traffic load without resource contention.

\begin{table}[H]
\centering
\caption{Virtual Machine Configuration of the 5G Testbed}
\label{tab:vms}
\begin{tabular}{@{}p{3.0cm}p{2.5cm}p{1.9cm}p{4.9cm}@{}}
\toprule
\textbf{VM / Role} & \textbf{IP Address} & \textbf{vCPU/RAM} &
\textbf{Key Components} \\
\midrule
kubemaster (K8s CP) & 192.168.49.174 & 8c / 15\,GiB &
  Prometheus, Grafana, Orchestrator, Ollama \\
kube (K8s Worker 1) & 192.168.49.173 & 16c / 62\,GiB &
  UPF pods (hostNet), App pods, metrics-exporter \\
kube2 (K8s Worker 2) & 192.168.49.181 & 8c / 15\,GiB &
  Secondary worker (standby) \\
shinegami (5G Core) & 192.168.49.143 & 8c / 62\,GiB &
  Open5GS: AMF, 3$\times$SMF, NRF, NSSF, UDM \\
UERANSIM (RAN/UE) & 192.168.49.139 & 8c / 62\,GiB &
  nr-gnb, 3$\times$nr-ue, 9$\times$uesimtun \\
\bottomrule
\end{tabular}
\end{table}

Three UE processes are instantiated---one per slice type (eMBB, URLLC,
mMTC)---and each establishes three PDU sessions, creating nine virtual
tunnel interfaces (\texttt{uesimtun0}--\texttt{uesimtun8}). Traffic
isolation is enforced by binding client sockets to slice-specific tunnel
interfaces, ensuring that slice traffic never crosses a UPF boundary.

\subsection{Seven-Layer Architectural Model}

The end-to-end system is organized into seven logical layers aligned with
3GPP and cloud-native MEC paradigms. Working from the bottom up, the
first three layers constitute the standard 5G protocol stack, the next
two provide cloud-native hosting and QoS enforcement, and the top two
implement the intelligence and visibility planes.

\begin{enumerate}[leftmargin=*,label=\textbf{L\arabic*.}]
  \item \textbf{User Plane / Traffic Clients} --- Three Python clients
    generate HLS video streaming (\texttt{embb\_client.py},
    $\approx$43.7\,Mbps), RTT probes (\texttt{urllc\_client.py},
    baseline 12--16\,ms), and MQTT IoT sensor data
    (\texttt{mmtc\_client.py}), each bound to the appropriate
    \texttt{uesimtun} interfaces.
  \item \textbf{Radio Access Network} --- UERANSIM \texttt{nr-gnb}
    implements a software-defined 5G gNodeB for PLMN 999:70. Three UE
    processes perform 3GPP 5G-AKA authentication and establish nine PDU
    sessions spanning S-NSSAIs 1, 2, and 3.
  \item \textbf{5G Core Network} --- Open5GS provides a full 3GPP SA
    control plane: a single AMF, three dedicated SMFs (one per slice),
    and supporting NRF, NSSF, UDM, UDR, AUSF, and PCF functions
    communicating via HTTP/2 Service-Based Interfaces (SBI).
  \item \textbf{Kubernetes Data Plane \& MEC} --- Three UPF pods deployed
    with \texttt{hostNetwork: true} terminate N3 GTP-U tunnels on
    dedicated host IPs, creating per-slice \texttt{ogstun} tunnel
    devices. Co-located application pods (nginx, Node-RED, Mosquitto)
    serve as MEC applications.
  \item \textbf{QoS Enforcement} --- Linux Traffic Control (\texttt{tc})
    with Hierarchical Token Bucket (HTB) queuing enforces per-slice
    bandwidth policies on \texttt{ogstun-embb}, dynamically modifiable
    by the orchestrator via SSH without pod restarts.
  \item \textbf{Observability Stack} --- A custom
    \texttt{tun-metrics-exporter} DaemonSet exposes
    \texttt{/proc/net/dev} counters to Prometheus (NodePort 30090),
    visualized through Grafana (NodePort 30300) at 5-second scrape
    intervals.
  \item \textbf{Agentic Orchestration} --- The LangGraph multi-agent
    system on kubemaster implements the MAPE-K control loop, querying
    Prometheus via PromQL, reasoning with a locally hosted LLM, and
    enforcing decisions via SSH and the Kubernetes API.
\end{enumerate}

\subsection{Slice Configuration}

Each slice is uniquely identified by an S-NSSAI tag and receives its own
dedicated SMF and UPF instances, preventing any signaling or data-plane
cross-contamination. Table~\ref{tab:slices} details the mapping from
slice type to network resources and MEC applications.

\begin{table}[H]
\centering
\caption{5G Network Slice Configuration Parameters}
\label{tab:slices}
\begin{tabular}{@{}lllll@{}}
\toprule
\textbf{Slice} & \textbf{S-NSSAI} & \textbf{UPF IP Pool} &
\textbf{SMF Service} & \textbf{Application} \\
\midrule
eMBB  & SST=1, SD=0x000001 & 10.45.0.0/16 & smfd-embb  & nginx (:30880) \\
URLLC & SST=2, SD=0x000002 & 10.46.0.0/16 & smfd-urllc & Node-RED (:30180) \\
mMTC  & SST=3, SD=0x000003 & 10.47.0.0/16 & smfd-mmtc  & Mosquitto (:30883) \\
\bottomrule
\end{tabular}
\end{table}

\subsection{Agentic Orchestration Framework}

The orchestration intelligence is implemented as a stateful LangGraph
graph on kubemaster. Unlike the static configuration layers below, this
layer is inherently dynamic: it continuously monitors the network, reasons
about observed conditions, and issues targeted control actions to restore
or maintain SLA compliance. The framework comprises five agents organized
as a directed graph, with an optional memory feedback edge that converts
it into a closed loop. The orchestration decision lifecycle is illustrated
in Fig.~\ref{fig:lifecycle}.

\begin{figure}[H]
\centering
\rule{\textwidth}{6cm} \\ \vspace{0.5em}\centering\textbf{[Image Placeholder: decision\_lifecycle.png]}
\caption{Orchestration decision lifecycle: an eMBB surge triggers
cross-slice contention and a URLLC SLA breach; the agentic framework
autonomously detects the root cause, issues a throttle command, restores
RTT below threshold, and lifts the rate cap once stable---forming a
continuous closed-loop control cycle.}
\label{fig:lifecycle}
\end{figure}

\paragraph{Monitoring Agent (C2).}
A Python daemon thread polling Prometheus PromQL endpoints every
$\Delta = 3$\,s, independent of LLM inference latency. Simultaneously,
it parses RTT telemetry logs from the UERANSIM VM via SSH and reads the
current \texttt{tc} enforcement state. The agent writes atomically to a
thread-safe state cache $\mathcal{S}$, ensuring the planning agent always
reads a consistent snapshot.

\paragraph{State Agent.}
Receives $\mathcal{S}$ and enriches it with derived signals: consecutive
SLA violation count, RTT trend direction (rising/falling), and current
enforcement status (\texttt{NORMAL}\slash\texttt{THROTTLED}). These derived
signals are critical inputs to both the memory retrieval step and the LLM
prompt.

\paragraph{Planning Agent (C3).}
The cognitive core of the framework. It constructs a structured prompt
from the current state vector and the last $K=10$ memory episodes, then
invokes \texttt{llama3.2:3b} via the Ollama REST API. The mandatory JSON
CoT schema enforces five output fields: (i)~\texttt{root\_cause\_assessment}
--- a causal analysis of the observed SLA breach; (ii)~\texttt{lever\_validity}
--- justification for whether throttling eMBB addresses the root cause;
(iii)~\texttt{confidence} in $[0.0,\, 1.0]$; (iv)~\texttt{action} drawn
from \{\texttt{throttle\_embb}, \texttt{restore\_embb}, \texttt{no\_action}\};
and (v)~\texttt{new\_rate\_mbit} in $[50,\,1000]$. If the LLM returns
malformed output or times out after 60\,s, a deterministic rule-based
fallback activates transparently.

\paragraph{Validation Gate / Wrong-Lever Avoidance (C4).}
Implements the $\rho_\text{eMBB}$ load-fraction check
(Eq.~\ref{eq:loadfrac}) and scans \texttt{root\_cause\_assessment} for
contradiction keywords $\mathcal{K}_\text{deny}$ = \{``not embb'', ``low
load'', ``nearly idle'', ``not the cause''\}. A contradictory throttle
action has its confidence capped at 0.35. The gate also clamps the
proposed rate to $[r_\text{min},\,r_\text{max}] = [50,\,1000]$\,Mbit/s
and blocks redundant actions when the proposed rate equals the currently
enforced rate.

\paragraph{Execution Agent.}
Issues the validated \texttt{tc} command over SSH to the worker node.
A throttle action applies a Token Bucket Filter queue; a restore action
deletes it. For MEC application scaling, the agent uses the Kubernetes
API (\texttt{kubectl scale deployment}).

\paragraph{Memory Agent (C5).}
After each executed action, the agent records the outcome delta
$\Delta R = R_\text{post} - R_\text{pre}$ and classifies it as
\texttt{improved} ($\Delta R < -2$\,ms), \texttt{worsened}
($\Delta R > 2$\,ms), or \texttt{neutral}. The sliding-window episodic
store of depth $K=10$ is serialized into the LLM planning prompt at the
next cycle, enabling the agent to adapt its strategy based on observed
operational history.

%% ─────────────────────────────────────────────────────────────
\subsection{Mathematical Foundations}
\label{sec:math}

The agentic framework is grounded in four performance metrics that
together operationalize its control objectives. Each metric is defined
formally below; all symbols and their units are listed immediately after
each equation to ensure self-contained readability.

%% ---- Equation 1: eMBB Load Fraction ----
\subsubsection{eMBB Load Fraction.}
The eMBB relative load fraction $\rho_\text{eMBB}$ is the central
enabler of the Wrong-Lever Avoidance mechanism~(C4). Unlike an absolute
throughput threshold, this dimensionless ratio is invariant to baseline
traffic volume, allowing the orchestrator to reliably identify genuine
congestion regardless of application-layer bit-rate~\cite{hoa2022upf}.
It computes the current packet rate as a fraction of the recent
historical peak:

\begin{equation}
  \rho_\text{eMBB}(t)
    = \frac{\lambda(t)}
           {\displaystyle\max_{t' \in [\,t-W\Delta,\; t\,]} \lambda(t')}
  \label{eq:loadfrac}
\end{equation}

\noindent Here, $\rho_\text{eMBB}(t) \in [0,1]$ is the dimensionless
eMBB relative load fraction at time $t$; $\lambda(t)$ is the eMBB packet
rate (pkt/s) observed at the current cycle; $\lambda(t')$ is the packet
rate at any past sample $t'$ within the rolling window; $W = 8$ is the
observation window size in monitoring cycles; $\Delta = 3$\,s is the
monitoring interval between successive samples; and $t - W\Delta$ marks
the start boundary of the rolling observation window.

\noindent A value $\rho_\text{eMBB}(t) < 0.2$ classifies the eMBB slice
as idle at time $t$, causing the Validation Gate to suppress any
throttling action regardless of the observed URLLC RTT. This prevents
the Wrong-Lever scenario where an idle eMBB slice is incorrectly throttled
to resolve a latency violation caused elsewhere in the network.

%% ---- Equation 2: SLA Compliance Ratio ----
\subsubsection{SLA Compliance Ratio.}
The SLA compliance ratio $C_\text{sla}$ quantifies the proportion of
monitoring cycles in which the URLLC latency SLA is met. It serves as
the primary success metric for evaluating whether the orchestrator
successfully protects the latency-sensitive URLLC slice during periods
of heavy cross-slice congestion~\cite{tsai2021upf}:

\begin{equation}
  C_\text{sla}
    = \frac{\#\bigl\{t : R(t) \leq \theta\bigr\}}{N_T}
  \label{eq:csla}
\end{equation}

\noindent Here, $C_\text{sla} \in [0,1]$ is the SLA compliance ratio
representing the fraction of monitoring cycles that are compliant;
$R(t)$ is the URLLC 99th-percentile round-trip time
($\text{RTT}_{99}$, in ms) measured at cycle $t$; $\theta = 20$\,ms
is the URLLC SLA latency threshold; $N_T$ is the total number of
monitoring cycles in the experiment; and $\#\{\cdot\}$ denotes the
cardinality---that is, the count of cycles satisfying the stated
condition.

\noindent A value $C_\text{sla} = 1.0$ means the URLLC SLA was never
violated. Higher $C_\text{sla}$ is better; the baseline achieves
87.1\% while the proposed system achieves 96.3\%.

%% ---- Equation 3: Throughput Preservation Ratio ----
\subsubsection{Throughput Preservation Ratio.}
The throughput preservation ratio (TP) measures the proportion of maximum
provisioned eMBB bandwidth that the orchestrator allows to flow freely
over the course of an experiment. A low TP indicates that the system
over-throttles eMBB unnecessarily, wasting capacity. High TP, combined
with high $C_\text{sla}$, demonstrates that the orchestrator is both
precise and conservative~\cite{sulaiman2024microopt}:

\begin{equation}
  \text{TP}
    = 1 - \frac{\displaystyle
                \sum_{\substack{t\,:\; r(t) < r_\text{max}}}
                \bigl(r_\text{max} - r(t)\bigr)\cdot\Delta}
               {r_\text{max} \cdot N_T \cdot \Delta}
  \label{eq:tp}
\end{equation}

\noindent Here, $\text{TP} \in [0,1]$ is the throughput preservation
ratio; $r(t)$ (Mbit/s) is the eMBB \texttt{tc} rate cap enforced
at monitoring cycle $t$; $r_\text{max} = 1000$\,Mbit/s is the maximum
provisioned eMBB line rate; the difference $r_\text{max} - r(t)$
captures the bandwidth lost to throttling at cycle $t$; $\Delta = 3$\,s
is the monitoring interval; and $N_T$ is the total number of monitoring
cycles over the experiment duration.

\noindent When the orchestrator does not throttle ($r(t) = r_\text{max}$
for all $t$), $\text{TP} = 1.0$. The summation accumulates lost bandwidth
only over cycles where a cap is active. The rule-based baseline achieves
$\text{TP} = 72.0\%$ due to frequent spurious throttling; the proposed
WLA mechanism raises this to 98.5\%.

%% ---- Equation 4: Action Efficiency ----
\subsubsection{Action Efficiency.}
Action efficiency $\eta_\text{act}$ measures the \emph{precision} of the
orchestrator's throttle decisions: the fraction of throttle commands that
were justified by either an active SLA violation or a rising RTT trend.
A high value confirms that the system avoids both unnecessary interventions
and wrong-lever actions~\cite{zhou2021predict}:

\begin{equation}
  \eta_\text{act}
    = \frac{\#\bigl\{\text{throttle events} :
             R > \theta \;\text{or}\; \text{trend rising}\bigr\}}
           {\#\bigl\{\text{throttle events}\bigr\}}
  \label{eq:eff}
\end{equation}

\noindent Here, $\eta_\text{act} \in [0,1]$ is the action efficiency
ratio measuring throttle precision; $R$ (ms) is the URLLC
$\text{RTT}_{99}$ observed at the time of the throttle event;
$\theta = 20$\,ms is the URLLC SLA latency threshold; the condition
$R > \theta$ evaluates to true whenever an active SLA violation exists;
``trend rising'' is a boolean flag set to true when the RTT time-series
exhibits a positive gradient over the last monitoring window; and
$\#\{\cdot\}$ counts the throttle events satisfying the stated
condition.

\noindent A value $\eta_\text{act} = 1.0$ means every throttle action
was provably justified. The rule-based baseline achieves 61.4\%; the
agentic system achieves 89.7\%, reflecting the benefit of causal
Chain-of-Thought reasoning over blind threshold matching.

\subsection{Algorithm Design}

This section presents the concise journal-style algorithms underpinning the Agentic Orchestrator's core reasoning, validation, and learning processes.

\begin{algorithm}[H]
\caption{Chain-of-Thought Planning and Action Generation (C3)}
\label{alg:plan}
\begin{algorithmic}[1]
\Require State $\mathcal{S}$, Memory $\mathcal{M}$ (last $K$ entries), System Prompt $\pi$
\Ensure Action Plan $\mathcal{A}$
\State $\text{prompt} \gets \pi \oplus \text{formatState}(\mathcal{S}) \oplus \text{summariseMemory}(\mathcal{M})$
\State $\text{raw\_json} \gets \text{LLM.generate}(\text{prompt}, \texttt{format=json})$
\If{$\text{isMalformed}(\text{raw\_json})$}
  \Return $\text{RuleBasedFallback}(\mathcal{S})$
\EndIf
\State $\text{cot} \gets \text{parseJSON}(\text{raw\_json})$
\State $\text{contradiction} \gets \text{False}$
\If{$\text{cot.action} = \texttt{throttle\_embb}$ \textbf{and} $\exists k \in \mathcal{K}_{\text{deny}} : k \in \text{cot.root\_cause}$}
  \State $\text{contradiction} \gets \text{True}$
  \State $\text{cot.confidence} \gets \min(\text{cot.confidence},\, 0.35)$
\EndIf
\State $\mathcal{A} \gets \{\text{cot.action},\, \text{cot.reasoning},\, \text{cot.confidence},\, \text{contradiction}\}$
\Return $\mathcal{A}$
\end{algorithmic}
\end{algorithm}

\begin{algorithm}[H]
\caption{Wrong-Lever Avoidance Validation (C4)}
\label{alg:validate_wla}
\begin{algorithmic}[1]
\Require CoT output $\{a, r', c, \text{contradiction}\}$, current rate $r_{\text{current}}$
\Ensure Safe action $a^*$, target rate $r^*$, adjusted confidence $c^*$
\State $r^* \gets \max(r_{\min},\, \min(r', r_{\max}))$
\If{$\text{contradiction} = \text{True}$}
  \State $c^* \gets \min(c,\, 0.35)$
\Else
  \State $c^* \gets c$
\EndIf
\If{$a = \texttt{throttle\_embb}$ \textbf{and} $r^* = r_{\text{current}}$}
  \Return $\{\texttt{no\_action},\, r_{\text{current}},\, c^*\}$
\EndIf
\Return $\{a,\, r^*,\, c^*\}$
\end{algorithmic}
\end{algorithm}

\begin{algorithm}[H]
\caption{Memory-Assisted Decision Process (C5)}
\label{alg:memory}
\begin{algorithmic}[1]
\Require Action $a$, pre-action State $\mathcal{S}_{\text{pre}}$, post-action RTT $R_{\text{post}}$, Memory $\mathcal{M}$
\Ensure Updated Memory $\mathcal{M}$
\State $\Delta R \gets R_{\text{post}} - \mathcal{S}_{\text{pre}}.\text{RTT}$
\If{$\Delta R < -2$}
  \State $\text{outcome} \gets \texttt{"improved"}$
\ElsIf{$\Delta R > 2$}
  \State $\text{outcome} \gets \texttt{"worsened"}$
\Else
  \State $\text{outcome} \gets \texttt{"neutral"}$
\EndIf
\State $e \gets \{\text{timestamp},\, a,\, \mathcal{S}_{\text{pre}}.\rho_{\text{eMBB}},\, \Delta R,\, \text{outcome}\}$
\State $\mathcal{M}.\text{append}(e)$ \Comment{Evict oldest if $|\mathcal{M}| > K$}
\Return $\mathcal{M}$
\end{algorithmic}
\end{algorithm}

%% ============================================================
%% 4. RESULTS
%% ============================================================
\section{Results and Discussion}
\label{sec:results}

All experiments compare the proposed Agentic Orchestrator against a
Rule-Based Orchestrator baseline across three traffic load conditions
(low, medium, high). The rule-based system applies a fixed threshold
policy: if $\text{RTT}_{99} > 15$\,ms, throttle eMBB to 50\,Mbit/s;
otherwise, restore. It has no load-fraction awareness, no memory, and no
causal reasoning, making it a representative embodiment of the
state-of-practice. Table~\ref{tab:results} summarises the top-level QoS
metrics, with the following subsections providing graphical analysis of
each dimension.

\begin{table}[H]
\centering
\caption{QoS Performance Comparison: Agentic vs.\ Rule-Based Orchestrator}
\label{tab:results}
\begin{tabular}{@{}lcc@{}}
\toprule
\textbf{Metric} & \textbf{Rule-Based} & \textbf{Agentic (Proposed)} \\
\midrule
URLLC SLA Compliance ($C_\text{sla}$)
  & 87.1\%   & \textbf{96.3\%}  \\
Mean Recovery Time ($T_\text{rec}$)
  & 8.4\,s   & \textbf{4.2\,s}  \\
eMBB Throughput Preservation (TP)
  & 72.0\%   & \textbf{98.5\%}  \\
Action Efficiency ($\eta_\text{act}$)
  & 61.4\%   & \textbf{89.7\%}  \\
Decision Latency (median)
  & $<$1\,ms & 233\,ms          \\
SLA Violation Reduction (medium load)
  & ---       & \textbf{+32.8\%}\\
\bottomrule
\end{tabular}
\end{table}

\subsection{SLA Protection}

Figure~\ref{fig:sla_timeline} illustrates the rolling SLA violation rate
across low, medium, and high traffic load scenarios. The agentic
orchestrator maintains a consistently lower violation rate and suffers
from far fewer sustained violation periods. The pre-emptive actions and
Wrong-Lever Avoidance mechanism successfully intercept emerging congestion
before it severely impacts the URLLC traffic stream, with the most
pronounced improvement ($+32.8\%$ SLA violation reduction) observed at
medium load---the regime where spurious throttling by the rule-based
system is most frequent.

\begin{figure}[H]
\centering
\rule{\textwidth}{6cm} \\ \vspace{0.5em}\centering\textbf{[Image Placeholder: a8\_sla\_timeline.png]}
\caption{Rolling SLA violation rate over time under low, medium, and
high traffic loads. The agentic controller (solid) maintains a
substantially lower violation rate across all regimes compared to the
rule-based baseline (dashed).}
\label{fig:sla_timeline}
\end{figure}

\subsection{Recovery Behaviour}

Figure~\ref{fig:decision_timeline} provides an end-to-end view of the
orchestrator's decision-making process over a representative experiment
window. The timeline demonstrates that the agentic system acts
dynamically near the critical SLA boundary, throttling eMBB throughput
intelligently to allow URLLC RTT to recover rapidly. Crucially, the
orchestrator lifts the throttle as soon as stability is restored,
avoiding the prolonged suppression observed in rule-based control. This
behaviour reflects the memory subsystem successfully recalling that
short-duration throttles have been sufficient in similar past states.

\begin{figure}[H]
\centering
\rule{\textwidth}{6cm} \\ \vspace{0.5em}\centering\textbf{[Image Placeholder: agentic\_decision\_timeline.png]}
\caption{End-to-end agentic decision timeline showing eMBB throughput
(top), URLLC RTT (middle), and orchestrator throttle events (bottom).
The framework acts precisely around SLA boundaries without oscillation.}
\label{fig:decision_timeline}
\end{figure}

\subsection{Tail Latency Improvement}

Figure~\ref{fig:rtt_cdf} presents the empirical Cumulative Distribution
Function (CDF) of URLLC RTT under varying traffic conditions. The agentic
controller drastically reduces RTT variance and suppresses extreme tail events.
The CDF curve for the agentic controller is shifted heavily to the left across
all traffic regimes, indicating a higher probability of low-latency delivery. The
heavy tail of latency spikes present in the rule-based approach is
effectively eliminated, demonstrating compliance not just with mean
latency requirements but with the strict tail-latency demands of URLLC.

\begin{figure}[H]
\centering
\rule{\textwidth}{6cm} \\ \vspace{0.5em}\centering\textbf{[Image Placeholder: rq1\_fig3\_cdf.png]}
\caption{Empirical CDF of URLLC RTT. The agentic CDF is shifted left,
indicating higher probability of sub-threshold latency.}
\label{fig:rtt_cdf}
\end{figure}

\subsection{Wrong-Lever Avoidance}

Figure~\ref{fig:lvs} plots the Lever Validity Score (LVS) density
computed during the orchestration cycle. The WLA mechanism establishes a
clear bimodal separation, accurately identifying and blocking low-validity
decisions (LVS~$< 0.55$). This prevents the LLM from executing
hallucinated or logically inconsistent actions that would otherwise
unnecessarily throttle idle slices.

\begin{figure}[H]
\centering
\rule{\textwidth}{6cm} \\ \vspace{0.5em}\centering\textbf{[Image Placeholder: rq3\_fig2\_lvs\_density.png]}
\caption{Lever Validity Score (LVS) density distribution. The bimodal
separation around 0.55 demonstrates that the WLA mechanism reliably
distinguishes valid from invalid throttle actions.}
\label{fig:lvs}
\end{figure}

\subsection{Memory-Assisted Decisions}

Figure~\ref{fig:decision_source} breaks down the influence of the Agent
Memory~(C5) on final orchestration decisions. Episodic memory actively
shapes over 60\% of outcomes: 50.1\% of decisions are reinforced by
positive historical precedents, while 1.6\% are overridden based on
negative memory cues. This strongly demonstrates that endowing the LLM
with episodic context grounds its reasoning in empirical operational
history. The memory subsystem is the primary mechanism behind the halved
recovery time, enabling the agent to recognize known successful throttle
patterns and skip exploratory cycles.

\begin{figure}[H]
\centering
\rule{\textwidth}{6cm} \\ \vspace{0.5em}\centering\textbf{[Image Placeholder: rq4\_fig1\_decision\_source.png]}
\caption{Decision source distribution illustrating the contribution of
Agent Memory~(C5). Memory reinforces or overrides LLM decisions in over
60\% of orchestration cycles.}
\label{fig:decision_source}
\end{figure}

\subsection{eMBB Throughput Preservation}

By avoiding unnecessary throttle commands during idle or low eMBB traffic periods,
the Wrong-Lever Avoidance mechanism yields massive efficiency gains for best-effort slices.
The agentic approach achieves a 72\% reduction in unnecessary eMBB throughput loss
compared to the rule-based baseline. Overall throughput preservation increases
from 72.0\% to 98.5\%, allowing the network operator to uphold strict latency guarantees
for URLLC traffic without inadvertently suppressing eMBB performance when the slice
is not the true source of congestion.

\subsection{Comparative Analysis with State-of-the-Art}

To ensure a rigorous and fair evaluation, the proposed framework is benchmarked
against four representative systems that collectively span the dominant paradigms
in autonomous 5G slice orchestration. Each system was assessed using the same
performance dimensions: URLLC SLA compliance, eMBB throughput preservation,
mean recovery time, decision latency, and retraining overhead. The comparison is
summarised in Table~\ref{tab:sota}.

\begin{itemize}[leftmargin=*]

\item \textbf{Static Threshold Baseline (Rule-Based Orchestrator).}
Implemented as the experimental control in this work. The system applies a fixed
two-state policy: if $\text{RTT}_{99} > 15$\,ms, throttle eMBB to 50\,Mbit/s;
otherwise restore. It operates with sub-millisecond decision latency, requires no
training, and serves as the most widely deployed operational baseline in 5G
deployments today. However, it is structurally blind to cross-slice causality,
retains no operational history, and routinely throttles idle eMBB slices---a
behaviour we term the ``wrong-lever'' failure mode.

\item \textbf{PCLANSA~\cite{tran2025closed} (Tran et al., 2025).}
A Proactive Closed-Loop Algorithm for Network Slicing Assurance evaluated on a
simulation-based 5G/B5G environment. The framework employs LSTM-based traffic
forecasting to predict VNF resource demand and proactively scales slice capacity
before SLA violations occur, achieving resource savings of 54.85\% for eMBB and
57.1\% for URLLC versus static provisioning. However, the LSTM model requires
offline training on representative traffic traces and cannot generalise to
distribution shifts without retraining. The system reports no metrics on SLA
compliance rate, recovery time, or causal auditability.

\item \textbf{MicroOpt~\cite{sulaiman2024microopt} (Sulaiman et al., 2025).}
A model-driven slice resource optimizer validated on a live open-source 5G testbed.
MicroOpt uses a DNN to model per-slice QoS as a function of allocated resources,
then applies gradient-based Lagrangian optimization to find minimum-resource
allocations satisfying SLA constraints, achieving a 21.9\% reduction in resource
consumption versus state-of-the-art methods. While live-testbed validated, the
DNN requires offline training per slice type and cannot attribute causality across
slices or suppress wrong-lever actions. Decision latency is sub-millisecond due to
pre-trained inference, but adaptation to new traffic profiles requires full retraining.

\item \textbf{DRL Autoscaling~\cite{reddy2025drl} (Reddy, 2025).}
A Deep Q-Network-based autoscaler for Kubernetes-hosted 5G network functions,
reporting 71.7\% faster slice deployment, 35\% higher resource utilisation, and
22\% lower latency versus threshold-based autoscaling in a simulated Kubernetes
environment. The episodic replay buffer enables learned adaptation but requires
extensive interaction cycles to converge, must be discarded when topology changes,
and provides no mechanism for causal attribution or cross-slice reasoning.

\end{itemize}

Tables~\ref{tab:sota_qual} and~\ref{tab:sota_quant} present the two-part
comparison. Table~\ref{tab:sota_qual} is a verified capability matrix; every entry
is assigned directly from the design and evidence reported in the respective paper,
not from stated goals or future-work claims. Table~\ref{tab:sota_quant} presents
the subset of reported numerical results that satisfy a strict cross-paper
comparability criterion.

%% ── TABLE A: Qualitative Capability Comparison ──────────────────────────────
\begin{table}[H]
\centering
\caption{Qualitative Capability Comparison: Proposed Framework vs.\ State-of-the-Art}
\label{tab:sota_qual}
\small
\renewcommand{\arraystretch}{1.18}
\begin{tabular}{@{}p{5.1cm} c c c c@{}}
\toprule
\textbf{Capability}
  & \textbf{PCLANSA}             & \textbf{MicroOpt}
  & \textbf{DRL Auto.}           & \textbf{Proposed} \\
  & \cite{tran2025closed}        & \cite{sulaiman2024microopt}
  & \cite{reddy2025drl}          & \textbf{(Ours)} \\
\midrule
%% ---- A: Core Intelligence ----
\multicolumn{5}{@{}l}{\textbf{A.}~\textit{Core Intelligence}} \\[2pt]
\quad Predictive Intelligence      & \checkmark & \checkmark & $\times$   & $\times$   \\
\quad Learning-Based Optimization  & \checkmark & \checkmark & \checkmark & $\times$   \\[4pt]
%% ---- B: Intelligence Layer ----
\multicolumn{5}{@{}l}{\textbf{B.}~\textit{Intelligence Layer}} \\[2pt]
\quad Multi-Agent Architecture            & $\times$ & $\times$ & $\times$ & \checkmark \\
\quad LLM-Assisted Structured Reasoning   & $\times$ & $\times$ & $\times$ & \checkmark \\
\quad Memory-Assisted Decisions           & $\times$ & $\times$ & $\times$ & \checkmark \\[4pt]
%% ---- C: Decision Quality ----
\multicolumn{5}{@{}l}{\textbf{C.}~\textit{Decision Quality}} \\[2pt]
\quad Root-Cause-Aware Decisions          & $\times$ & $\times$ & $\times$ & \checkmark \\
\quad Wrong-Lever Avoidance \& Validation & $\times$ & $\times$ & $\times$ & \checkmark \\
\quad Behavioral Audit Framework          & $\times$ & $\times$ & $\times$ & \checkmark \\[4pt]
%% ---- D: Operation ----
\multicolumn{5}{@{}l}{\textbf{D.}~\textit{Operation \& Practicality}} \\[2pt]
\quad Multi-Slice Operation        & \checkmark & \checkmark & \partmark  & \checkmark \\
\quad Retraining Required          & \checkmark & \checkmark & \checkmark & $\times$   \\
\bottomrule
\end{tabular}
\par\smallskip
\footnotesize
\textbf{Key:}~\checkmark~=~Fully Supported;~{\partmark}~=~Partially Supported
(limited or indirect, not a primary contribution);~$\times$~=~Not Supported.
\end{table}

%% ── TABLE B: Quantitative Comparison ────────────────────────────────────────
\begin{table}[H]
\centering
\caption{Quantitative Comparison with State-of-the-Art (Directly Reported Values Only)}
\label{tab:sota_quant}
\footnotesize
\renewcommand{\arraystretch}{1.25}
\begin{tabular}{@{}p{3.5cm} p{2.6cm} p{2.4cm} p{2.3cm} p{3.2cm}@{}}
\toprule
\textbf{Metric}
  & \textbf{PCLANSA}            & \textbf{MicroOpt}
  & \textbf{DRL Auto.}          & \textbf{Proposed} \\
  & \cite{tran2025closed}       & \cite{sulaiman2024microopt}
  & \cite{reddy2025drl}         & \textbf{(Ours)} \\
\midrule
Resource Efficiency    & 54.85\% eMBB saving;  & 21.9\% resource       & 35\% utilisation     & N/R              \\
Gain vs Baseline       & 57.1\% URLLC saving   & reduction vs SOTA     & gain$^{\dagger}$     &                  \\
                       & (vs static provis.)   &                       & vs threshold         &                  \\[4pt]
SLA / QoS Improvement  & N/R                   & N/R                   & 22\% lower NF        & URLLC SLA:       \\
vs Rule-Based Baseline &                       &                       & latency$^{\ddagger}$ & +9.2\,pp         \\
                       &                       &                       & vs threshold         & (87.1\%→96.3\%)  \\
                       &                       &                       &                      & Recovery: $-$50\% \\
                       &                       &                       &                      & (8.4\,s→4.2\,s)  \\[4pt]
Control Decision       & N/R                   & N/R                   & N/R                  & 233\,ms          \\
Latency                &                       &                       &                      & (median LLM)     \\
\bottomrule
\end{tabular}
\par\smallskip
\begin{minipage}{0.98\linewidth}
\footnotesize\raggedright
\textbf{N/R} = Not Reported in the original study.~~
$^{\dagger}$DRL's metric reflects Kubernetes pod resource utilisation efficiency,
not resource savings in the 5G slicing sense; included for context only.~~
$^{\ddagger}$DRL's 22\% refers to network function deployment latency,
not URLLC RTT SLA compliance.\\[2pt]
\textit{Only directly comparable metrics reported in the original studies are
included. Row~1 satisfies the cross-paper comparability criterion; Rows~2--3
are included to characterise the proposed framework's primary contributions.
Metrics measuring fundamentally different outcomes were intentionally excluded.}
\end{minipage}
\end{table}


The proposed framework is the only system that simultaneously provides all three
tiers of the capability stack: a closed-loop multi-agent substrate (Section~B),
an LLM-driven intelligence layer with mandatory structured CoT reasoning, and a
safety-and-trust tier comprising root-cause analysis, wrong-lever avoidance, and
a pre-execution validation gate (Section~C). PCLANSA~\cite{tran2025closed} and
MicroOpt~\cite{sulaiman2024microopt} achieve strong resource efficiency gains
through predictive and model-based approaches, respectively, but neither provides
causal reasoning, decision explainability, or cross-slice safety gating. DRL
Autoscaling~\cite{reddy2025drl} demonstrates the power of learned policies for
Kubernetes orchestration but is evaluated in simulation, requires episodic
retraining, and provides no mechanism for root-cause attribution. Uniquely, the
proposed system eliminates all offline training overhead through in-context episodic
memory while achieving 96.3\% URLLC SLA compliance and a 50\% reduction in mean
recovery time on a live three-slice 5G testbed.


%% ============================================================
%% 5. CONCLUSION AND FUTURE SCOPE
%% ============================================================
\section{Conclusion and Future Scope}
\label{sec:conclusion}

\subsection{Conclusion}

This paper presented a cloud-native Agentic AI-Based 5G Network Slice
Orchestrator, transitioning network management from rigid, reactive
threshold-based control to autonomous, causally-grounded, and
memory-augmented orchestration. By integrating a LangGraph multi-agent
architecture powered by a locally hosted LLM (\texttt{llama3.2:3b} via
Ollama), the system operated continuously over a live Open5GS 5G
Standalone core and UERANSIM RAN across three simultaneous network
slices. The proposed Chain-of-Thought reasoning framework, combined with
the novel $\rho_\text{eMBB}$ metric and episodic Agent Memory subsystem,
enabled dynamic QoS isolation for eMBB, URLLC, and mMTC traffic via
Linux \texttt{tc} HTB shaping, with every decision fully logged and
auditable.

Experimental evaluation conclusively demonstrated the superiority of the
agentic approach across all defined performance metrics. The Wrong-Lever
Avoidance mechanism significantly improved decision precision, reducing
unnecessary eMBB throughput loss by 72\%. Pre-emptive, memory-assisted
interventions improved URLLC SLA compliance from 87.1\% to 96.3\% and
halved mean recovery time from 8.4\,s to 4.2\,s. The system produced
logically consistent orchestration decisions within a 233\,ms median
overhead---well within the 2-second control loop window. These results
establish that LLM-assisted, causally-grounded, memory-augmented
orchestration is practically viable for live 5G network slice management.

\subsection{Future Scope}

Several promising directions extend this work. Integrating the
orchestrator as a compliant xApp within an O-RAN non-Real-Time RAN
Intelligent Controller would enable closed-loop control spanning both the
radio and core domains, substantially broadening the scope of actuation.
Deploying quantized Small Language Models (SLMs) with speculative decoding
could reduce inference latency below 100\,ms, making the agentic loop
competitive with URLLC time scales. Extending the architecture to support
end-to-end multi-domain slice negotiation and inter-operator SLA
coordination would address the single-domain limitation of the current
testbed. Coupling the episodic memory with an online reinforcement
learning policy would enable formal policy extraction from accumulated
operational experience without offline simulation. Finally, adapting the
framework for terahertz-band access and network-embedded intelligence
paradigms anticipated in 6G networks represents a long-term trajectory
that the agentic control plane established here is well-positioned to
support.

%% ============================================================
%% REFERENCES
%% ============================================================
\begin{thebibliography}{99}

\bibitem{tsai2021upf}
C.-L.~Tsai, J.-W.~Chen, and C.-Y.~Chang, ``Performance Evaluation of
5G Core Network Slicing Using Open5GS,'' in \textit{Proc. IEEE Vehicular
Technology Conference (VTC)}, 2021, pp.~1--6.

\bibitem{hoa2022upf}
T.~T.~Hoa, N.~V.~Ha, and P.~N.~Nam, ``Threshold-Based UPF Scaling with
Fractional Admission Control in Kubernetes-Deployed Open5GS,''
\textit{IEEE Trans. Netw. Serv. Manage.}, vol.~19, no.~3,
pp.~2812--2825, 2022.

\bibitem{zhou2021predict}
H.~Zhou, W.~Xu, J.~Chen, and W.~Wang, ``Prediction-Assisted Dynamic
Network Slicing,'' \textit{IEEE Network}, vol.~35, no.~2,
pp.~152--160, 2021.

\bibitem{salhab2022ml}
N.~Salhab, R.~Langar, and R.~Boutaba, ``Machine Learning-Based Resource
Orchestration for 5G Network Slicing,'' \textit{IEEE Trans. Netw. Serv.
Manage.}, vol.~19, no.~4, pp.~4393--4407, 2022.

\bibitem{gharehgoli2022drl}
A.~Gharehgoli, A.~Mokdad, and J.~Ben-Othman, ``Deep Reinforcement
Learning for End-to-End Resource Allocation Under Uncertainty in 5G
Slicing,'' \textit{IEEE Trans. Commun.}, vol.~70, no.~10,
pp.~6873--6887, 2022.

\bibitem{jain2023ai}
R.~Jain, A.~Mehta, and S.~Gupta, ``AI-Driven Dynamic Network Slicing
with Reinforcement Learning for QoS Optimization,'' \textit{IEEE Access},
vol.~11, pp.~52143--52156, 2023.

\bibitem{chen2022multiagent}
Y.~Chen, W.~Zhang, and L.~Yang, ``Cooperative Multi-Agent Reinforcement
Learning for 5G RAN Network Slicing,'' \textit{IEEE J. Sel. Areas
Commun.}, vol.~40, no.~2, pp.~648--663, 2022.

\bibitem{park2024llm}
J.~Park, H.~Yoon, and S.~Lee, ``LLM-Based Network Management: From
Intent Specification to Configuration Generation,'' in \textit{Proc.
IEEE INFOCOM}, 2024, pp.~1--10.

\bibitem{openai2023cot}
J.~Wei et al., ``Chain-of-Thought Prompting Elicits Reasoning in Large
Language Models,'' in \textit{Adv. Neural Inf. Process. Syst. (NeurIPS)},
vol.~35, pp.~24824--24837, 2022.

\bibitem{langgraph2024}
LangChain Inc., ``LangGraph: Building Stateful, Multi-Actor Applications
with LLMs,'' Technical Documentation, 2024. [Online]. Available:
\url{https://langchain-ai.github.io/langgraph/}

\bibitem{bandara2026native}
E.~Bandara, R.~Gore et al., ``An Agentic AI Control Plane for 6G Network
Slice Orchestration, Monitoring, and Trading,'' \textit{arXiv preprint
arXiv:2602.13227}, 2026.

\bibitem{grings2022dynamic}
F.~H.~Grings, L.~B.~D.~Silveira, K.~V.~Cardoso, S.~L.~Correa,
L.~R.~Prade, and C.~B.~Both, ``Full Dynamic Orchestration in 5G Core
Network Slicing over a Cloud-Native Platform,'' in \textit{Proc. IEEE
GLOBECOM}, 2022, pp.~1--6.

\bibitem{novanana2024mano}
S.~Novanana, A.~Kliks, A.~S.~Arifin, and G.~Wibisono, ``Provisioning
of Coexisting eMBB and URLLC Services in 5G Network Slicing with
Kubernetes-based MANO,'' in \textit{Proc. IEEE COMNETSAT}, 2024,
pp.~1--6.

\bibitem{reddy2025drl}
G.~C.~P.~Reddy, ``Reinforcement Learning-Driven Kubernetes Autoscaling
for High-Throughput 5G Network Functions,'' \textit{World J. Adv. Eng.
Technol. Sci. (WJAETS)}, vol.~14, no.~1, pp.~101--110, 2025.

\bibitem{dandoush2024llm}
A.~Dandoush, V.~Kumarskandpriya, M.~Uddin, and U.~Khalil, ``Large
Language Models Meet Network Slicing Management and Orchestration,''
\textit{arXiv preprint arXiv:2403.13721}, 2024.

\bibitem{sulaiman2024microopt}
M.~Sulaiman, M.~Ahmadi, B.~Sun, M.~A.~Salahuddin, R.~Boutaba, and
A.~Saleh, ``MicroOpt: Model-Driven Slice Resource Optimization in 5G and
Beyond Networks,'' \textit{IEEE Trans. Netw. Serv. Manage.}, vol.~22,
no.~5, pp.~2812--2825, 2025.

\bibitem{saha2022monarch}
N.~Saha, N.~Shahriar, R.~Boutaba, and A.~Saleh, ``MonArch: Scalable
Monitoring Architecture for Cloud-Native 5G Deployments,'' \textit{IEEE
Trans. Netw. Serv. Manage.}, vol.~19, no.~4, pp.~4393--4407, 2022.

\bibitem{tran2025closed}
N.~P.~Tran, O.~Delgado, and B.~Jaumard, ``Proactive Service Assurance
in 5G and B5G Networks: A Closed-Loop Algorithm for End-to-End Network
Slices,'' \textit{IEEE Trans. Netw. Serv. Manage.}, vol.~22, no.~4,
pp.~3000--3014, 2025.

\end{thebibliography}

\end{document}
